% PART 3. CONCLUSION (Phần kết luận) - mục đánh số 1., 2., 1.1.
\addchap{PART 3. CONCLUSION}\label{chap:conclusion}
\ThesisFlatNumbering

\section{Provisional Findings at the Progress Cutoff}
\subsection{Knowledge and Skills}
The work has produced a traceable computational workflow from experimental read counts to mutation
labels, protein-group splits and model comparisons. The initial restricted-PSI studies established
ridge as a strong baseline and showed that extending training-label coverage can help even when a
learning curve within the original distribution is nearly flat. Negative findings on interaction
designs and tile-label pre-training were retained rather than replaced by unsupported architecture claims.

The full-PSI extension evaluates \nVthreeNPeptides{} peptides in \nVthreeNGroups{} protein groups.
Full-range W5 ridge training reduces macro MAE from \nVthreeRidgeB{} to \nVthreeRidgeD{} PSI.
Matched frozen-ESM controls provide development evidence for positional mixing: the CNN reaches
\nVthreeCnnD{}, compared with \nVthreePointwiseD{} for pointwise prediction. Scanning-label
pre-training helps its one-hot CNN control, whereas quality weighting does not help ridge.
These are conclusions about specified populations and training recipes, not a universal model ranking.

The completed local expansion adds all nine planned pLM checkpoints, 540 frozen-feature ridge fits
and 60 matched ESM2-35M frozen/LoRA fits. LoRA reduces MAE from \nCampaignFrozenB{} to
\nCampaignLoraB{} in arm B and from \nCampaignFrozenD{} to \nCampaignLoraD{} in arm D;
the adjusted paired intervals support both improvements. Larger frozen encoders do not consistently
improve performance, and the nested family selector reaches \nCampaignNested{} MAE. B and D
remain close overall while differing across WT-PSI strata. These results support a specified adaptation
recipe and a comparison procedure rather than a universal family ranking.

Out-of-fold three-class calibration is complete for the four matched core recipes. LoRA improves
Brier score, but positive-event ECE and high-probability reliability still expose errors. The locked
W5-D transfer evaluation is also complete: positive AUROC is \nCampaignTAuc{}, whereas the MAE
difference from predicting zero is \nCampaignTDelta{} with interval \nCampaignTCI{}.
The contrast between ranking and regression limits the transfer claim.

\subsection{Program and Practical Application Potential}
The calibrated restricted-domain ridge baseline has an inference path numerically consistent with
the research runner. A local React/FastAPI implementation supplies single-peptide results, reference
lookup, annotations and traceable exports. Its intended use is to shortlist substitutions for
experimental follow-up, with visible model, source and domain evidence.
The newer workspace UI demonstrates the intended interactions with synthetic fixtures; it is not
an integrated persistent application. Projects/history, durable peptide/protein jobs, experimental
import/comparison, PostgreSQL integration and their local acceptance remain unfinished as of
6 November 2026. The serving model has not been replaced by a development leaderboard result.

\section{Limitations and Challenges}
Predictive evidence comes from public reporter-assay sources that have informed several design
decisions, including a locked transfer subset from a previously explored source. Protein-group
splitting reduces detectable overlap but does not establish
complete biological independence; thousands of mutation cells do not create thousands of independent
peptides. The original 91-peptide evaluation and the full-PSI task have different populations and
metrics. Most follow-up comparisons are exploratory, and the seeds share an outer partition.

The executable local graph is complete; reviewed external-source and structure inputs remain pending.
The locked transfer set covers R/A/S substitutions with replicate-1 labels in the same laboratory.
High-WT-PSI regression and positive-probability shift show that favourable development ranking does
not establish robust transfer. LoRA has not been evaluated on T. Runtime, precision and feature
choices also differ across model families. CV calibrators are not final serving artifacts; model-specific
final refit, lineage cards and export remain unfinished. Earlier baseline calibration and its domain
gate cannot be assumed valid for a new serving model.
The gate is not yet integrated correctly into the baseline application. No new wet-lab validation,
E3-specific recognition, intact-protein prediction or clinical inference has been established.

\section{Priorities for the Next Research Phase}
Preserve the completed OOF record, paired uncertainty and unsuccessful controls. Investigate the
high-PSI transfer errors and probability shift with an explicitly designed new evaluation, rather than
using the already inspected T set to choose among pLMs. New external and structure extensions require
reviewed inputs and overlap/rights audits. Any model selected for inference needs an architecture-specific
final refit, calibration, domain policy, model card and export; the benefit of the three-seed ensemble
must be retained or re-evaluated if the serving configuration changes.

Complete the agreed local software workflows and check numerical consistency, correct domain
handling, coordinate/assay alignment, saved-work recovery and measured resource use. Update the
remaining installation and acceptance material from executed evidence. Further data sources or
biological tests require endpoint, homology and source-rights audits before any broader claim.
